\documentclass{article}

\usepackage{../header}

\title{Principles of Quantum Mechanics}

\author{Notes made by Finley Cooper}


\begin{document}

\maketitle

\newpage

\tableofcontents

\newpage

\section{Introduction}

Recall from IB Quantum Mechanics we have the following:

\begin{enumerate}

    \item Wave-particle duality. Waves behave like particles, like light quanta, and vice versa (particles behave like waves).

    \item We have a wavefunction $\psi(x)$ for a particle which we interpret as a probability density $|\psi(x)|^2$ (i.e. probability is intrinsic to the theory).

    \item Observables become operators on wavefunctions:

        \begin{itemize}
            \item $\hat{x} \to x$;
            \item $\hat{p} \to -i\hbar \frac{\partial}{\partial x}$
        \end{itemize}

        The lack of commutation, since $[\hat{x},\hat{p}] = i\hbar$, leads to the uncertainty relation.

    \item Schrödinger's equation with Hamiltonian $H$ specifies dynamics, and eigenvalues of $H$ are the allowed energy levels of the system.

    \item All of this explains the hydrogen atom structure and the resulting line spectra.

\end{enumerate}

The aims of this course are the following:

\begin{enumerate}

    \item Reformulate quantum mechanics in a more powerful, flexible and useful form: the Dirac formulation, which goes beyond wavefunctions. The previous description appears as a particular realisation of our generalisation of quantum states and also leads to alternatives such as

        \[
            \tilde{\psi}(p)
            \quad \text{with} \quad
            \hat{x} \to i\hbar \frac{\partial}{\partial p}
            \quad \text{and} \quad
            \hat{p} \to p.
        \]

    \item This more general approach allows much simpler analysis of known problems (for example, the Simple Harmonic Oscillator), but also lots of new things such as \textit{spin} of particles and symmetries (translations, rotations) and the related conservation laws.

    \item We will also look at identical particles (bosons and fermions).

    \item We will develop a framework for quantising more general systems such as EM fields.

    \item We will also set up perturbation theory, which is a practical tool for computation, but we will not heavily consider the detailed applications.

\end{enumerate}

\section{Dirac Formalism}

\subsection{Concepts and Definitions}

A quantum system is described at each instant by a \textit{state} $\ket{\psi}$ which belongs to a complex vector space $V$; any non-zero element of this space $V$ is a possible state. Thus

\[
    \ket{\psi}, \ket{\phi} \in V
    \implies
    \alpha\ket{\psi} + \beta\ket{\phi} \in V
    \quad \text{for any } \alpha,\beta \in \mathbb{C}.
\]

Physically, this is a superposition principle leading to wave-like behaviour, but these states are \textit{not wavefunctions}. There are also \textit{dual} or conjugate states $\bra{\psi}$ which belong to the dual space $V^\dagger$. By definition, states and duals can be combined or paired to give a complex number. The map

\[
    (\bra{\psi}, \ket{\phi}) \mapsto \bra{\psi}\ket{\phi}
\]

defines a map from $V^\dagger \times V \to \mathbb{C}$ with

\[
    \bra{\phi}
    \left(\alpha_1\ket{\psi_1}+\alpha_2\ket{\psi_2}\right)
    =
    \alpha_1\bra{\phi}\ket{\psi_1}
    +
    \alpha_2\bra{\phi}\ket{\psi_2}
\]

and

\[
    \left(\beta_1\bra{\phi_1}+\beta_2\bra{\phi_2}\right)\ket{\psi}
    =
    \beta_1\bra{\phi_1}\ket{\psi}
    +
    \beta_2\bra{\phi_2}\ket{\psi}.
\]

The spaces $V$ and $V^\dagger$ come with an inner product which can be described as a one-to-one correspondence between the spaces

\begin{align*}
    V &\to V^\dagger \\
    \ket{\psi} &\to \bra{\psi}
    = \left(\ket{\psi}\right)^\dagger.
\end{align*}

We also have that if we take a linear combination of states to the dual space, the coefficients have to have their complex conjugates taken when mapping from the dual space to $V$ or vice versa. Hence the inner product is

\[
    \ket{\phi},\ket{\psi}
    \mapsto
    \bra{\phi}\ket{\psi}
    =
    \left(\ket{\phi}\right)^\dagger\ket{\psi}
\]

and it is assumed to obey

\[
    \bra{\phi}\ket{\psi}
    =
    \left(\bra{\psi}\ket{\phi}\right)^*
\]

since the inner product is Hermitian, and the inner product defines a positive-definite norm on $V$:

\[
    \|\ket{\psi}\|^2 = \bra{\psi}\ket{\psi}.
\]

\begin{remark}

    Knowing $\bra{\phi}\ket{\psi}$ for all $\bra{\phi}$ determines $\ket{\psi}$ and vice versa.

\end{remark}

We can show this by

\[
    \bra{\phi}\ket{\psi_1}
    =
    \bra{\phi}\ket{\psi_2}
    \implies
    \bra{\phi}\ket{\chi}=0,
\]

where $\ket{\chi}=\ket{\psi_1}-\ket{\psi_2}$. If this holds for all $\bra{\phi}$, we can set $\bra{\phi}=\bra{\chi}$, so by positive definiteness we have that the norm of $\ket{\chi}$ is zero, and therefore $\ket{\chi}=0$ too.

\par

We also have that the physical content of any state is unchanged by

\[
    \ket{\psi} \to \alpha\ket{\psi}
\]

for any non-zero $\alpha \in \mathbb{C}$. So we will usually normalise to make the state have norm $1$, but we can still change it by a phase factor.

\begin{remark}

    Relative phases can be important in expressions where we're adding a linear combination of elements of $V$. We also assume the space $V$ is complete so appropriate sequences and series converge in $V$.

\end{remark}

A complete inner product space is called a \textit{Hilbert space}.

\par

\begin{definition}

    (Operator) An \textit{operator} $Q$ is a linear map on states

    \[
        \ket{\psi} \to Q\ket{\psi}.
    \]

\end{definition}

The same operator can be thought of as acting to the left on dual states, so

\[
    \bra{\phi} \to \bra{\phi}Q
\]

where the right-hand side is defined by

\[
    \left(\bra{\phi}Q\right)\ket{\psi}
    =
    \bra{\phi}\left(Q\ket{\psi}\right)
    =
    \bra{\phi}Q\ket{\psi}.
\]
\begin{remark}
	States and operators $ \ket \psi, \bra \phi, Q $ can be dealt with by rulers that mirror those for complex vectors $ \mathbf u (\text{col}), \mathbf v^\dagger(\text{row}) $ and a matrix $ A $. So we can compare $ \bra \psi Q \ket \phi $ with $ \mathbf{v}^\dagger A\mathbf u $ like we do with inner products.
\end{remark}
\begin{definition}
	(Eigenstate) For any operator $ Q $ with say that $ \ket \psi $ non-zero is an \textit{eigenstate} with \textit{eigenvalue} $ \lambda $ if 
	\[
	  Q\ket \psi = \lambda \ket \psi
	\]
	or equivalently
	\[
	  \bra \psi Q^\dagger = \lambda^*\bra \psi.
	\]
\end{definition}
\begin{remark}
  For a general $ Q $, $ \lambda $ may be complex.
\end{remark}
\begin{definition}
	(Hermitian Conjugate) For any operator $ Q $, the \textit{Hermitian conjugate} or \textit{adjoint map} $ Q^\dagger $ is an operator defined by
	\[
	  \bra \psi Q^\dagger = (Q\ket \psi)^\dagger
	\]
	for any $ \ket \psi $.
\end{definition}
Equivalently
\begin{align*}
	\bra \phi Q^\dagger \ket \psi &= (Q\bra \phi)^\dagger \ket \psi \\
				      &= \bra \psi Q \ket \phi^\dagger
\end{align*}
for all $ \ket\psi,\ket \phi $.
\par
\begin{definition}
	(Hermitian) An operator $ Q $ is \textit{Hermitian} or \textit{self-adjoint} if
	\[
		Q^\dagger = Q
	\]
	and such operators are called \textit{observables} because they correspond to physical quantities that can be measured, for example position, momentum, energy.
\end{definition}
To define a quantum system, we must specify the \textit{Hamiltonian} whose eigenvalues are the allowed energy levels of the system,
\[
  H\ket \chi = E\ket \chi.
\]
We have the following key results for any Hermitian operator $ Q $.
\begin{proposition}
  For a Hermitian operator $ Q $,
  \begin{enumerate}
	  \item All eigenvalues are real;
	  \item Eigenstates with distinct eigenvalues are orthogonal;
	  \item Any state can be written as a linear combination of its eigenstates (i.e its matrix representation is diagonalisable).
  \end{enumerate}
\end{proposition}
  For a proof of (iii) we have a proof in IB Linear Algebra. We'll cover a proof of (i) and (ii) here.\par
  \pf Suppose that $ Q\ket \psi = \lambda\ket \psi $. So taking Hermitian conjugates on both sides we have that 
  \begin{align*}
	  \bra \psi Q^\dagger &= \lambda^*\bra \psi \\
			      &= \bra \psi Q.
  \end{align*}
  Hence $ \bra \psi Q \ket \psi = \lambda \bra{\psi}\ket{\psi} = \lambda^* \bra{\psi}\ket{\psi} $. Since $ \ket \psi $ is non-zero, it's norm is non-zero hence $ \lambda = \lambda^* $.

\par
For the second part let $ \ket n $ be an eigenstate of $ Q $ with eigenvalues $ \lambda = q_n $ real. Then we have that
\begin{align*}
	Q \ket n &= q_n\ket n \\
	Q\ket m &= q_m\ket m \\
	\bra m Q&= q_m\bra m
\end{align*}
hence $ \bra m Q \ket n = q_n\bra{m}\ket{n} = q_m\bra{m}\ket{n} $. But we have that $ q_m\ne q_n $ so $ \bra{m}\ket{n} = 0 $.\qed
\subsection{The Harmonic Oscillator}
\subsubsection{Construction of the eigenstates}
A one-dimensional harmonic oscillator of mass $ m $ and frequency $ \omega $ is a system with position $ \hat x $, momentum $ \hat p $ and Hamilitonian
\[
	 H = \frac 1{2m} \hat p^2 + \frac  12 m \omega^2 \hat x^2.
\]
We seek the energy eigenstates and eigenvalues of $ H $. We do this by introducting \textit{annihilation, creation, and number operators}.\par
Define
\begin{align*}
	a &= \left(\frac {m\omega}{2\hbar}\right)^{\frac 12}\left(\hat x + \frac{i\hat  p }{m \omega}\right) \\
	a^\dagger &= \left(\frac {m\omega}{2\hbar}\right)^{\frac 12}\left(\hat x - \frac{i\hat p}{m\omega}\right).
	\end{align*}
Equivalently 
\begin{align*}
	\hat x &= \left(\frac \hbar{2 m \omega}\right)^{\frac 12} \left( a + a^\dagger\right) \\
   \hat p &= \left(\frac {\hbar m \omega}2\right)^{\frac 12}i\left( a^\dagger - a\right).
\end{align*}
It's easy to check that
\[
	[\hat x, \hat p] = i\hbar \iff [a, a^\dagger] = 1.
\]
In fact,
\begin{align*}
	aa^\dagger  &= \frac{m\omega}{2\hbar} \hat x^2 + \frac 1{2m\omega \hbar} \hat p ^2 - \frac i{2\hbar} (\hat x \hat p - \hat p \hat x) \\
	&= \frac 1{\hbar \omega} H + \frac 12
\end{align*}
but
\[
	a^\dagger a = \frac 1{\hbar \omega} H - \frac 12.
\]
This confirms the commutation relation, but we also find that
\begin{align*}
  H = \hbar \omega \left(N + \frac 12\right)
\end{align*}
with $ N = a^\dagger a $, the \textit{number operator} which is clearly Hermitian, so finding eigenvalues and eigenstates of $ N $ is equivalent to finding these for $ H $.\par
Let $ \ket \lambda $ be any eigenstate of $ N $ with eigenvalue $ \lambda $, so
\[
  N \ket \lambda = \lambda \ket \lambda \iff H \ket \lambda = E \ket \lambda
\]
where $ E = \hbar \omega\left(\lambda + \frac 12 \right) $. Then assuming that $ \ket \lambda $ is normalised, write
\[
	\lambda = \bra\lambda N \ket \lambda = \lambda a^\dagger a \ket \lambda = (a\ket\lambda)^\dagger (a\ket \lambda) = ||a\ket \lambda||^2
\]
we could have just seen this directly, since $ N $ is clearly positive semi-definite from its construction as the outer product of an operator. So we have that $ \lambda \ge 0  $ and $ \lambda = 0 \iff a\ket \lambda = 0 $.\par 
Now we consider
\begin{align*}
	N(a^\dagger \ket \lambda) & = ([N, a^\dagger] + a^\dagger N)\ket \lambda \\
				  &= ( a^\dagger + a^\dagger \lambda) \ket \lambda \\
				  &= (\lambda + 1) (a^\dagger \ket \lambda).
\end{align*}
Similarly we have that
\begin{align*}
	N(a\ket \lambda) = &= ([N,a] + aN)\ket \lambda\\
			   &= ( -a + a\lambda)\ket\lambda \\
			   &= (\lambda -1 ) a \ket \lambda.
\end{align*}
Here we have used the results for commutators $ [N,a^\dagger] = [a^\dagger, a^\dagger] = a^\dagger[a,a^\dagger] = a^\dagger $ and similarily $ [N, a] = -a $.
\par
Now the news states with be eigenstates provided their norms are non-zero. We can compute 
\begin{align*}
	||a\ket \lambda ||^2 &= \lambda\\
	||a^\dagger \ket \lambda ||^2 &= \bra \lambda aa^\dagger \ket \lambda = \lambda a^\dagger a + 1 \ket \lambda \\
			     &= \lambda + 1.
\end{align*}
Because of these properties $ a^\dagger $ and $ a $ are called \textit{creation} (or raising) and \textit{annihilation} (or lowering) operators.\par
We can't lower below zero, since the eigenvalues are positive so the eigenvalues must be a subset of the positive integers to avoid a contradiction. We can lower and raise all we want, so we can generate any other state with eigenvalues being all of the non-negative integers.\par
More formally, since we know that $ \lambda \ge 0  $, so for a given eigenstate $ \ket \lambda $, if $ \lambda $ is \textit{not} an integer, then
\[
  a\ket \lambda, a^2\ket\lambda, \dots, a^n\ket \lambda
\]
are states with eigenvalues
\[
  \lambda -1, \lambda -2, \dots, \lambda -n
\]
but for $ n $ sufficiently large $ \lambda - n <0 $ which is a contradiction.
\par
By contrast if $ \lambda =k \in \N $ then the contradiction is avoided because
\[
  a\ket k, a^2 \ket k, \dots, a^k \ket k
\]
with eigenvalues
\[
  k-1, k-2, \dots, 0
\]
so we can't lower anymore because our eigenstate has collapsed to zero, i.e. if we force integer eigenvalues the resultant eigenstate is \textit{nilpotent}.\par
We can use the raising operator $ a^\dagger $ to generate the eigenstates with eigenvalues $ k+1, k+2, \dots $ forever.
\par
Hence the eigenvalues of $ N $ are therefore precisely the non-negative integers and the oscillator has energy levels
\[
  E_n = \hbar \omega\left(n+\frac 12 \right).
\]
From our calculation of norms above, we can choose normalised eigenstates $ \ket n $ related by 
\begin{align*}
	a^\dagger \ket n &= \sqrt{n+1}\ket{n+1} \\
	a\ket n &= \sqrt{n} \ket{n-1}.
\end{align*}
Starting from any state, $ \ket n $, we can generate all others. In particular, starting from the \textit{ground state} $ \ket 0 $, characteristed by 
\[
  a\ket 0 = 0
\]
we can construct normalised eigenstates
\[
	\ket n =\frac 1{\sqrt{n!}} (a^\dagger)^n \ket 0.
\]
\begin{remark}
	Beware that $ \ket 0 $ is \textit{not} zero in $ V $. It is an eigenstate and therefore non-zero by definition.
\end{remark}
\subsubsection{Importance of the oscillator}
\begin{enumerate}
	\item A smooth potential $ V(x) $ with equilibrium point $ x = x_0 $ (so that $ V'(x_0) = 0 $) can be expanded by
		\[
		  V(x_0 + x) = V(x_0) + \frac 12 V''(x_0) x^2 + O(x^3)
		\]
		so, if $ x=x_0 $ is \textit{stable}, then we can set $ V''(x_0) = m\omega^2 > 0 $ and for displacements $ x $ not too large, we can expect the system to be well approximated by a oscillator.
	\item This can also be applied in the case for a diatomic molecules where we have quantisation of vibrational energies which is important in understanding the internal energy and heat capacities which result in macroscopic consequences.
	\item More complicated systems can be analysed in terms of normal modes. Coherent motion in which all degrees of freedom oscillate with the same common frequency $ \omega $. The general solution for a classical oscillator is
		\[
			x(t) = A e^{-i\omega t} + A^* e^{i\omega t}
		\]
		where $ A $ is a complex constant.\par
		Normal modes for a system with variables $ x_N $ are solutions of the form
		\[
			x_N(t) = Au_Ne^{-i\omega t} + A^* u^*_N e^{i\omega t}
		\]
		with $ u_N $ a complex vector specific to each normal frequency $ \omega $.
	\item For example consider a crystal with a number of atoms $ \mathcal N \sim 10^{23} $ where the force are approximately elastic and $ N $ labels diretions of motion for each atom, so we have $ 3\times \mathcal N $ values.
	\item Or we consider an EM field, where normal modes are oscillations of electric and mangetic at each point in space are given by
		\[
			\mathbf E(\x,t) = A\mathbf u(\x) e^{-i\omega t} + A^*\mathbf u(\x)^*e^{i\omega t}.
		\]
		This is exact.
	\item Normal modes are independent oscillator: treat these as above. For each mode we have
		\[
			a(\omega) \quad\text{and}\quad a(\omega)^\dagger
		\]
		where are ground state $ \ket 0 $ with $ a(\omega)\ket 0  = 0$. Non-ground states created with the raising operators add energy and momentum and behave like particles. In this way we can see that $ a $ destroys particles and $ a^\dagger $ creates particles. 
\end{enumerate}
\subsection{Bases and Measurements}
\subsubsection{Bases of Eigenstates}
We know that for any observable $ Q $ there is an orthonormal basis of eigenstates $ \{\ket n\} $ for $ V $,
\[
  Q \ket n = q_n\ket n.
\]
We have orthonormality $ \bra{m}\ket{n} = \delta_{mn} $ and for any state $ \ket \psi \in V $ we have that
\[
  \ket \psi = \sum_n \alpha_n \ket n
\]
where $ \alpha_n = \bra{n}\ket{\psi} $. Note that $ ||\ket \psi||^2 = 1 $ if and only if
\[
	\bra{\psi}\ket{\psi} = \left(\sum_m \alpha_m^*\bra m\right)\left(\sum_n\alpha_n \ket n \right) = \sum_n |\alpha_n|^2 = 1.
\]
Note that we can write
\[
  1 = \sum_n \ket n \bra n.
\]
since for any state $ \ket \psi $ we have that
\begin{align*}
	\ket \psi &= \left(\sum_n \ket n \bra n\right)\ket \psi \\
		  &= \sum_n \ket n(\bra{n}\ket{\psi})
\end{align*}
In above there may be several independent staets with the same eigenvalue $ \lambda $. The eigenspace for a given eigenvalue is
\[
	V_\lambda = \{ \ket \psi : Q\ket \psi = \lambda \ket \psi\}.
\]
Note that $ V_\lambda $ has basis $ \{ \ket n : q_n = \lambda\} $. The \textit{degeneracy} of $ \lambda $ is the number of sattes in this basis (i.e the dimension of $ V_\lambda $). If $ V_\lambda $ has dimension $ 1 $, then we say that $ \lambda $ is non-degenerate.












\end{document}
